\documentclass[12pt,a4paper]{article}
\usepackage[utf8]{vietnam}
\usepackage{amsmath,amssymb}
\usepackage{enumerate}
\usepackage{tabularx}
\usepackage{tikz}
\usepackage{xcolor}
\usepackage[top=0.45cm, bottom=0.45cm, left=0.8cm, right=0.8cm]{geometry}
\DeclareUnicodeCharacter{00B2}{\ensuremath{^2}}
\DeclareUnicodeCharacter{00B3}{\ensuremath{^3}}

\setlength{\parskip}{0pt}
\setlength{\parindent}{0pt}
\setlength{\fboxsep}{3.5pt}
\setlength{\fboxrule}{0.5pt}

% Dòng kẻ chấm mờ cự ly chuẩn viết tay tiểu học (18pt)
\newcommand{\dongke}[1]{%
  \begingroup%
  \setlength{\parskip}{0pt}%
  \baselineskip=18pt%
  \vspace{0.4mm}%
  \foreach \i in {1,...,#1}{%
    \noindent{\color{gray!55}\dotfill}\par%
  }%
  \vspace{0.2mm}%
  \endgroup%
}

% Khung bao quanh bài tập tự luận
\newcommand{\khungBaiTap}[1]{%
  \vspace{0.8mm}%
  \noindent\fbox{%
    \parbox{\dimexpr\linewidth-2\fboxsep-2\fboxrule\relax}{%
      \vspace{0.3mm}%
      #1%
      \vspace{0.3mm}%
    }%
  }%
  \vspace{0.5mm}%
}

% Ô bài tập đặt tính
\newcommand{\oDatTinh}[3][4]{%
  \vspace{0.3mm}%
  \noindent\textbf{#2} #3\par%
  \dongke{#1}%
}

% Câu trắc nghiệm dàn 4 đáp án trên 1 hàng ngang
\newcommand{\cauTN}[6]{%
  \noindent\textbf{Câu #1.} #2\par\vspace{0.1mm}%
  \noindent%
  \makebox[0.25\linewidth][l]{\textbf{A.} #3}%
  \makebox[0.25\linewidth][l]{\textbf{B.} #4}%
  \makebox[0.25\linewidth][l]{\textbf{C.} #5}%
  \makebox[0.25\linewidth][l]{\textbf{D.} #6}\par\vspace{0.6mm}%
}

% Câu trắc nghiệm dàn 2 đáp án / hàng
\newcommand{\cauTNhaiDong}[6]{%
  \noindent\textbf{Câu #1.} #2\par\vspace{0.1mm}%
  \noindent%
  \makebox[0.5\linewidth][l]{\textbf{A.} #3}%
  \makebox[0.5\linewidth][l]{\textbf{B.} #4}\par\vspace{0.1mm}%
  \noindent%
  \makebox[0.5\linewidth][l]{\textbf{C.} #5}%
  \makebox[0.5\linewidth][l]{\textbf{D.} #6}\par\vspace{0.6mm}%
}

\begin{document}
\pagestyle{empty}
\fontsize{10.2pt}{13.2pt}\selectfont

% =============================================================================
% ========================= MÃ ĐỀ: 401 ========================================
% =============================================================================

% -------------------- TRANG 1: MẶT TRƯỚC --------------------
\noindent
\begin{tabularx}{\linewidth}{|X|c|}
\hline
\textbf{PHIẾU BÀI TẬP CUỐI TUẦN -- TOÁN LỚP 4 (TUẦN 12)} & \textbf{MÃ ĐỀ: 401} \\
\textit{Trung tâm Toán Thầy Lĩnh -- Năm học: 2026 - 2027} & \textit{Thời gian: 45 phút} \\
\hline
\end{tabularx}

\vspace{1mm}

\noindent
\begin{tabularx}{\linewidth}{|X|p{3.5cm}|}
\hline
\rule[-2mm]{0pt}{7.5mm}\textbf{Họ và tên:} \dotfill & \textbf{Lớp:} \makebox[2cm]{\dotfill} \\
\hline
\end{tabularx}

\vspace{1mm}

\noindent
\begin{tabularx}{\linewidth}{|p{3cm}|X|}
\hline
\centering\textbf{ĐIỂM SỐ} & \centering\textbf{LỜI NHẬN XÉT CỦA THẦY CÔ} \tabularnewline
\hline
\rule{0pt}{1.1cm} & \tabularnewline
\hline
\end{tabularx}

\vspace{1mm}

\noindent\textbf{BẢNG ĐIỀN ĐÁP ÁN TRẮC NGHIỆM (Mỗi câu đúng 0,25 điểm):}
\vspace{0.6mm}

\begin{tabularx}{\linewidth}{|c|*{12}{>{\centering\arraybackslash}X|}}
\hline
\textbf{Câu} & \textbf{1} & \textbf{2} & \textbf{3} & \textbf{4} & \textbf{5} & \textbf{6} & \textbf{7} & \textbf{8} & \textbf{9} & \textbf{10} & \textbf{11} & \textbf{12} \\
\hline
\textbf{Đáp án} & & & & & & & & & & & & \\
\hline
\end{tabularx}

\vspace{1.2mm}

\noindent\textbf{\large PHẦN I. TRẮC NGHIỆM KHÁCH QUAN (3,0 điểm)}

\cauTN{1}{Kết quả của phép tính $654\,820 - 238\,475$ là:}{$416\,345$}{$416\,445$}{$426\,345$}{$415\,345$}
\cauTN{2}{Biểu thức $(a + b) + c = a + (b + c)$ thể hiện tính chất nào?}{Giao hoán phép cộng}{Kết hợp phép cộng}{Phân phối phép nhân}{Giao hoán phép nhân}
\cauTN{3}{Tính thuận tiện: $45 + 137 + 55 = \dots$}{$227$}{$237$}{$247$}{$257$}
\cauTN{4}{Số thích hợp điền vào chỗ chấm $6\,\text{tấn}\,450\,\text{kg} = \dots\,\text{kg}$ là:}{$645$}{$6\,450$}{$60\,450$}{$64\,500$}
\cauTN{5}{Số thích hợp điền vào chỗ chấm $9\,\text{m}^2\,35\,\text{dm}^2 = \dots\,\text{dm}^2$ là:}{$935$}{$9\,035$}{$9\,350$}{$98$}
\cauTN{6}{Một thế kỉ rưỡi ($1\,\text{thế kỉ} + \dfrac{1}{2}\,\text{thế kỉ}$) bằng bao nhiêu năm?}{$125$ năm}{$150$ năm}{$175$ năm}{$200$ năm}
\cauTN{7}{Tìm số bị trừ: $x - 34\,280 = 51\,640$:}{$x = 85\,920$}{$x = 17\,360$}{$x = 85\,820$}{$x = 86\,020$}
\cauTN{8}{Điền dấu thích hợp: $75\,000 - 25\,000 \;\dots\; 80\,000 - 32\,000$}{$>$}{$<$}{$=$}{Không so sánh được}
\cauTN{9}{Hiệu giữa số lớn nhất có 5 chữ số và số bé nhất có 5 chữ số là:}{$89\,999$}{$90\,000$}{$89\,990$}{$99\,998$}
\cauTN{10}{Làm tròn số $685\,400$ đến hàng chục nghìn ta được số:}{$680\,000$}{$690\,000$}{$685\,000$}{$700\,000$}
\cauTNhaiDong{11}{Một hình tam giác có độ dài ba cạnh lần lượt là $25\,\text{cm}, 35\,\text{cm}, 40\,\text{cm}$. Chu vi hình tam giác đó là:}{$90\,\text{cm}$}{$100\,\text{cm}$}{$110\,\text{cm}$}{$120\,\text{cm}$}
\cauTNhaiDong{12}{Một kho chứa $5\,000\,\text{kg}$ thóc, đã xuất đi $1\,\text{tấn}\,500\,\text{kg}$. Trong kho còn lại số ki-lô-gam thóc là:}{$3\,500\,\text{kg}$}{$4\,500\,\text{kg}$}{$3\,000\,\text{kg}$}{$2\,500\,\text{kg}$}

\vspace{0.4mm}
\noindent\textbf{\large PHẦN II. TỰ LUẬN (7,0 điểm)}

\noindent\textbf{Bài 1 (2,0 điểm). Đặt tính rồi tính:}
\vspace{0.4mm}

\noindent
\begin{tabularx}{\linewidth}{|X|X|X|X|}
\hline
\oDatTinh[4]{a)}{$742\,530 - 285\,648$} & 
\oDatTinh[4]{b)}{$850\,000 - 364\,725$} & 
\oDatTinh[4]{c)}{$458\,362 + 276\,845$} & 
\oDatTinh[4]{d)}{$614\,258 + 185\,742$} \tabularnewline
\hline
\end{tabularx}

\newpage
% -------------------- TRANG 2: MẶT SAU --------------------
\noindent
\begin{tabularx}{\linewidth}{|X|c|}
\hline
\textbf{PHẦN TỰ LUẬN (TIẾP THEO) -- PHIẾU BÀI TẬP TUẦN 12} & \textbf{MÃ ĐỀ: 401} \\
\hline
\end{tabularx}

\vspace{0.8mm}

% BÀI 2: TÌM X (4 Ý)
\khungBaiTap{%
  \textbf{Bài 2 (1,5 điểm). Tìm $x$:}\par\vspace{0.6mm}
  \noindent
  \begin{minipage}[t]{0.485\linewidth}
    \textbf{a)} $x + 38\,450 = 85\,120$\par
    \dongke{3}
    \vspace{0.4mm}
    \textbf{c)} $92\,500 - x = 38\,640$\par
    \dongke{3}
  \end{minipage}\hfill
  \begin{minipage}[t]{0.485\linewidth}
    \textbf{b)} $x - 47\,280 = 65\,390$\par
    \dongke{3}
    \vspace{0.4mm}
    \textbf{d)} $x - (15\,400 + 12\,600) = 45\,000$\par
    \dongke{3}
  \end{minipage}
}

\vspace{0.8mm}

% BÀI 3: ĐỔI ĐƠN VỊ (4 Ý - 1 CỘT TỪ TRÊN XUỐNG)
\khungBaiTap{%
  \textbf{Bài 3 (1,0 điểm). Viết số thích hợp vào chỗ chấm:}\par\vspace{0.8mm}
  \noindent
  a) $7\,\text{tấn}\,280\,\text{kg} = \dotfill\,\text{kg}$\par\vspace{1.4mm}
  \noindent
  b) $8\,\text{m}^2\,46\,\text{dm}^2 = \dotfill\,\text{dm}^2$\par\vspace{1.4mm}
  \noindent
  c) $4\,\text{thế kỉ}\,45\,\text{năm} = \dotfill\,\text{năm}$\par\vspace{1.4mm}
  \noindent
  d) $3\,\text{phút}\,40\,\text{giây} = \dotfill\,\text{giây}$\par
}

\vspace{0.8mm}

% BÀI 4: BÀI TOÁN GIẢI (11 DÒNG KẺ)
\khungBaiTap{%
  \textbf{Bài 4 (1,5 điểm).} Một cửa hàng nhập về 12{,}5 tạ gạo. Ngày thứ nhất cửa hàng bán được 40\% số gạo đó, ngày thứ hai bán được 65\% số gạo còn lại. Mỗi ki-lô-gam gạo cửa hàng bán với giá 18 000 đồng. Hỏi sau hai ngày, cửa hàng thu được tất cả bao nhiêu tiền bán gạo?
  \par\vspace{0.3mm}%
  \noindent\textit{Bài giải:}
  \dongke{11}
}

\vspace{0.8mm}

% BÀI 5: TOÁN NÂNG CAO (8 DÒNG KẺ)
\khungBaiTap{%
  \textbf{Bài 5 (1,0 điểm).} Tính bằng cách thuận tiện nhất:\par\vspace{0.6mm}\centerline{$K = 12{,}5 \times 101 - 12{,}5$}
  \par\vspace{0.3mm}%
  \noindent\textit{Bài làm:}
  \dongke{8}
}

\newpage


% =============================================================================
% ========================= MÃ ĐỀ: 402 ========================================
% =============================================================================

% -------------------- TRANG 1: MẶT TRƯỚC --------------------
\noindent
\begin{tabularx}{\linewidth}{|X|c|}
\hline
\textbf{PHIẾU BÀI TẬP CUỐI TUẦN -- TOÁN LỚP 4 (TUẦN 12)} & \textbf{MÃ ĐỀ: 402} \\
\textit{Trung tâm Toán Thầy Lĩnh -- Năm học: 2026 - 2027} & \textit{Thời gian: 45 phút} \\
\hline
\end{tabularx}

\vspace{1mm}

\noindent
\begin{tabularx}{\linewidth}{|X|p{3.5cm}|}
\hline
\rule[-2mm]{0pt}{7.5mm}\textbf{Họ và tên:} \dotfill & \textbf{Lớp:} \makebox[2cm]{\dotfill} \\
\hline
\end{tabularx}

\vspace{1mm}

\noindent
\begin{tabularx}{\linewidth}{|p{3cm}|X|}
\hline
\centering\textbf{ĐIỂM SỐ} & \centering\textbf{LỜI NHẬN XÉT CỦA THẦY CÔ} \tabularnewline
\hline
\rule{0pt}{1.1cm} & \tabularnewline
\hline
\end{tabularx}

\vspace{1mm}

\noindent\textbf{BẢNG ĐIỀN ĐÁP ÁN TRẮC NGHIỆM (Mỗi câu đúng 0,25 điểm):}
\vspace{0.6mm}

\begin{tabularx}{\linewidth}{|c|*{12}{>{\centering\arraybackslash}X|}}
\hline
\textbf{Câu} & \textbf{1} & \textbf{2} & \textbf{3} & \textbf{4} & \textbf{5} & \textbf{6} & \textbf{7} & \textbf{8} & \textbf{9} & \textbf{10} & \textbf{11} & \textbf{12} \\
\hline
\textbf{Đáp án} & & & & & & & & & & & & \\
\hline
\end{tabularx}

\vspace{1.2mm}

\noindent\textbf{\large PHẦN I. TRẮC NGHIỆM KHÁCH QUAN (3,0 điểm)}

\cauTN{1}{Kết quả của phép tính $742\,650 - 358\,284$ là:}{$384\,366$}{$384\,466$}{$394\,366$}{$384\,266$}
\cauTN{2}{Biểu thức $(28 + 19) + 81 = 28 + (19 + 81)$ là áp dụng tính chất:}{Giao hoán phép cộng}{Kết hợp phép cộng}{Giao hoán phép nhân}{Phân phối}
\cauTN{3}{Tính thuận tiện: $128 + 250 + 72 = \dots$}{$440$}{$450$}{$460$}{$470$}
\cauTN{4}{Số thích hợp điền vào chỗ chấm $8\,\text{tấn}\,320\,\text{kg} = \dots\,\text{kg}$ là:}{$832$}{$8\,320$}{$83\,200$}{$80\,320$}
\cauTN{5}{Số thích hợp điền vào chỗ chấm $6\,\text{m}^2\,58\,\text{dm}^2 = \dots\,\text{dm}^2$ là:}{$658$}{$6\,058$}{$6\,580$}{$69$}
\cauTN{6}{Khoảng thời gian $2\,\text{thế kỉ}$ bằng bao nhiêu năm?}{$100$ năm}{$200$ năm}{$300$ năm}{$2\,000$ năm}
\cauTN{7}{Tìm số trừ: $84\,500 - x = 36\,280$:}{$x = 48\,220$}{$x = 48\,320$}{$x = 120\,780$}{$x = 47\,220$}
\cauTN{8}{Điền dấu thích hợp: $85\,000 - 35\,000 \;\dots\; 90\,000 - 42\,000$}{$>$}{$<$}{$=$}{Không so sánh được}
\cauTN{9}{Hiệu giữa số lớn nhất có 6 chữ số và số lớn nhất có 5 chữ số là:}{$900\,000$}{$899\,999$}{$990\,000$}{$900\,001$}
\cauTN{10}{Làm tròn số $726\,800$ đến hàng chục nghìn ta được số:}{$720\,000$}{$730\,000$}{$726\,000$}{$700\,000$}
\cauTNhaiDong{11}{Một tam giác có chu vi $60\,\text{cm}$, biết hai cạnh dài $18\,\text{cm}$ và $22\,\text{cm}$. Độ dài cạnh còn lại là:}{$20\,\text{cm}$}{$25\,\text{cm}$}{$15\,\text{cm}$}{$30\,\text{cm}$}
\cauTNhaiDong{12}{Một bồn chứa $8\,000\,\text{lít}$ nước, người ta đã tháo ra $3\,400\,\text{lít}$. Trong bồn còn lại số lít nước là:}{$4\,600\,\text{lít}$}{$5\,600\,\text{lít}$}{$4\,400\,\text{lít}$}{$5\,400\,\text{lít}$}

\vspace{0.4mm}
\noindent\textbf{\large PHẦN II. TỰ LUẬN (7,0 điểm)}

\noindent\textbf{Bài 1 (2,0 điểm). Đặt tính rồi tính:}
\vspace{0.4mm}

\noindent
\begin{tabularx}{\linewidth}{|X|X|X|X|}
\hline
\oDatTinh[4]{a)}{$825\,460 - 378\,592$} & 
\oDatTinh[4]{b)}{$900\,000 - 452\,836$} & 
\oDatTinh[4]{c)}{$516\,428 + 294\,785$} & 
\oDatTinh[4]{d)}{$735\,614 + 164\,386$} \tabularnewline
\hline
\end{tabularx}

\newpage
% -------------------- TRANG 2: MẶT SAU --------------------
\noindent
\begin{tabularx}{\linewidth}{|X|c|}
\hline
\textbf{PHẦN TỰ LUẬN (TIẾP THEO) -- PHIẾU BÀI TẬP TUẦN 12} & \textbf{MÃ ĐỀ: 402} \\
\hline
\end{tabularx}

\vspace{0.8mm}

% BÀI 2: TÌM X (4 Ý)
\khungBaiTap{%
  \textbf{Bài 2 (1,5 điểm). Tìm $x$:}\par\vspace{0.6mm}
  \noindent
  \begin{minipage}[t]{0.485\linewidth}
    \textbf{a)} $x + 42\,680 = 91\,340$\par
    \dongke{3}
    \vspace{0.4mm}
    \textbf{c)} $95\,200 - x = 46\,370$\par
    \dongke{3}
  \end{minipage}\hfill
  \begin{minipage}[t]{0.485\linewidth}
    \textbf{b)} $x - 53\,490 = 74\,820$\par
    \dongke{3}
    \vspace{0.4mm}
    \textbf{d)} $x - (18\,200 + 13\,800) = 52\,000$\par
    \dongke{3}
  \end{minipage}
}

\vspace{0.8mm}

% BÀI 3: ĐỔI ĐƠN VỊ (4 Ý - 1 CỘT TỪ TRÊN XUỐNG)
\khungBaiTap{%
  \textbf{Bài 3 (1,0 điểm). Viết số thích hợp vào chỗ chấm:}\par\vspace{0.8mm}
  \noindent
  a) $9\,\text{tấn}\,160\,\text{kg} = \dotfill\,\text{kg}$\par\vspace{1.4mm}
  \noindent
  b) $7\,\text{m}^2\,32\,\text{dm}^2 = \dotfill\,\text{dm}^2$\par\vspace{1.4mm}
  \noindent
  c) $5\,\text{thế kỉ}\,35\,\text{năm} = \dotfill\,\text{năm}$\par\vspace{1.4mm}
  \noindent
  d) $4\,\text{phút}\,25\,\text{giây} = \dotfill\,\text{giây}$\par
}

\vspace{0.8mm}

% BÀI 4: BÀI TOÁN GIẢI (11 DÒNG KẺ)
\khungBaiTap{%
  \textbf{Bài 4 (1,5 điểm).} Một bể nước hình hộp chữ nhật có chiều dài 2{,}5 m, chiều rộng 1{,}6 m và chiều cao 1{,}2 m. Người ta dùng một máy bơm có công suất 0{,}8 $\text{m}^3$ nước mỗi giờ để bơm nước vào bể. Hỏi sau 2 giờ, máy bơm được bao nhiêu mét khối nước và lượng nước trong bể chiếm bao nhiêu phần trăm thể tích bể?
  \par\vspace{0.3mm}%
  \noindent\textit{Bài giải:}
  \dongke{11}
}

\vspace{0.8mm}

% BÀI 5: TOÁN NÂNG CAO (8 DÒNG KẺ)
\khungBaiTap{%
  \textbf{Bài 5 (1,0 điểm).} Tính bằng cách thuận tiện nhất:\par\vspace{0.6mm}\centerline{$K = 8{,}4 \times 9{,}9 + 8{,}4 \times 0{,}1$}
  \par\vspace{0.3mm}%
  \noindent\textit{Bài làm:}
  \dongke{8}
}

\newpage


% =============================================================================
% ========================= MÃ ĐỀ: 403 ========================================
% =============================================================================

% -------------------- TRANG 1: MẶT TRƯỚC --------------------
\noindent
\begin{tabularx}{\linewidth}{|X|c|}
\hline
\textbf{PHIẾU BÀI TẬP CUỐI TUẦN -- TOÁN LỚP 4 (TUẦN 12)} & \textbf{MÃ ĐỀ: 403} \\
\textit{Trung tâm Toán Thầy Lĩnh -- Năm học: 2026 - 2027} & \textit{Thời gian: 45 phút} \\
\hline
\end{tabularx}

\vspace{1mm}

\noindent
\begin{tabularx}{\linewidth}{|X|p{3.5cm}|}
\hline
\rule[-2mm]{0pt}{7.5mm}\textbf{Họ và tên:} \dotfill & \textbf{Lớp:} \makebox[2cm]{\dotfill} \\
\hline
\end{tabularx}

\vspace{1mm}

\noindent
\begin{tabularx}{\linewidth}{|p{3cm}|X|}
\hline
\centering\textbf{ĐIỂM SỐ} & \centering\textbf{LỜI NHẬN XÉT CỦA THẦY CÔ} \tabularnewline
\hline
\rule{0pt}{1.1cm} & \tabularnewline
\hline
\end{tabularx}

\vspace{1mm}

\noindent\textbf{BẢNG ĐIỀN ĐÁP ÁN TRẮC NGHIỆM (Mỗi câu đúng 0,25 điểm):}
\vspace{0.6mm}

\begin{tabularx}{\linewidth}{|c|*{12}{>{\centering\arraybackslash}X|}}
\hline
\textbf{Câu} & \textbf{1} & \textbf{2} & \textbf{3} & \textbf{4} & \textbf{5} & \textbf{6} & \textbf{7} & \textbf{8} & \textbf{9} & \textbf{10} & \textbf{11} & \textbf{12} \\
\hline
\textbf{Đáp án} & & & & & & & & & & & & \\
\hline
\end{tabularx}

\vspace{1.2mm}

\noindent\textbf{\large PHẦN I. TRẮC NGHIỆM KHÁCH QUAN (3,0 điểm)}

\cauTN{1}{Kết quả của phép tính $816\,430 - 452\,875$ là:}{$363\,555$}{$363\,655$}{$373\,555$}{$362\,555$}
\cauTN{2}{Cho phép tính $(a + b) + c$. Khi thay đổi vị trí dấu ngoặc thành $a + (b + c)$ thì kết quả:}{Không đổi}{Tăng lên}{Giảm đi}{Không tính được}
\cauTN{3}{Tính thuận tiện: $340 + 185 + 60 = \dots$}{$575$}{$585$}{$595$}{$600$}
\cauTN{4}{Số thích hợp điền vào chỗ chấm $5\,\text{tấn}\,620\,\text{kg} = \dots\,\text{kg}$ là:}{$562$}{$5\,620$}{$50\,620$}{$56\,200$}
\cauTN{5}{Số thích hợp điền vào chỗ chấm $4\,\text{m}^2\,72\,\text{dm}^2 = \dots\,\text{dm}^2$ là:}{$472$}{$4\,072$}{$4\,720$}{$49$}
\cauTN{6}{Khoảng thời gian $2\,\text{phút}\,18\,\text{giây}$ bằng bao nhiêu giây?}{$128\,\text{giây}$}{$138\,\text{giây}$}{$148\,\text{giây}$}{$158\,\text{giây}$}
\cauTN{7}{Tìm $x$: $x - 45\,380 = 62\,150$:}{$x = 107\,530$}{$x = 16\,770$}{$x = 107\,430$}{$x = 108\,530$}
\cauTN{8}{Điền dấu thích hợp: $92\,000 - 42\,000 \;\dots\; 85\,000 - 30\,000$}{$>$}{$<$}{$=$}{Không so sánh được}
\cauTN{9}{Hiệu giữa số bé nhất có 6 chữ số và số lớn nhất có 4 chữ số là:}{$90\,001$}{$90\,000$}{$89\,999$}{$90\,010$}
\cauTN{10}{Làm tròn số $843\,600$ đến hàng chục nghìn ta được số:}{$840\,000$}{$850\,000$}{$843\,000$}{$800\,000$}
\cauTNhaiDong{11}{Một tứ giác có độ dài các cạnh là $12\,\text{cm}, 15\,\text{cm}, 18\,\text{cm}, 25\,\text{cm}$. Chu vi tứ giác đó là:}{$60\,\text{cm}$}{$70\,\text{cm}$}{$65\,\text{cm}$}{$80\,\text{cm}$}
\cauTNhaiDong{12}{Một xưởng may có $6\,500\,\text{m}$ vải, đã dùng hết $2\,\text{km}\,300\,\text{m}$ vải. Xưởng còn lại số mét vải là:}{$4\,200\,\text{m}$}{$4\,300\,\text{m}$}{$3\,200\,\text{m}$}{$4\,500\,\text{m}$}

\vspace{0.4mm}
\noindent\textbf{\large PHẦN II. TỰ LUẬN (7,0 điểm)}

\noindent\textbf{Bài 1 (2,0 điểm). Đặt tính rồi tính:}
\vspace{0.4mm}

\noindent
\begin{tabularx}{\linewidth}{|X|X|X|X|}
\hline
\oDatTinh[4]{a)}{$934\,510 - 478\,635$} & 
\oDatTinh[4]{b)}{$700\,000 - 325\,486$} & 
\oDatTinh[4]{c)}{$482\,639 + 358\,742$} & 
\oDatTinh[4]{d)}{$568\,234 + 231\,766$} \tabularnewline
\hline
\end{tabularx}

\newpage
% -------------------- TRANG 2: MẶT SAU --------------------
\noindent
\begin{tabularx}{\linewidth}{|X|c|}
\hline
\textbf{PHẦN TỰ LUẬN (TIẾP THEO) -- PHIẾU BÀI TẬP TUẦN 12} & \textbf{MÃ ĐỀ: 403} \\
\hline
\end{tabularx}

\vspace{0.8mm}

% BÀI 2: TÌM X (4 Ý)
\khungBaiTap{%
  \textbf{Bài 2 (1,5 điểm). Tìm $x$:}\par\vspace{0.6mm}
  \noindent
  \begin{minipage}[t]{0.485\linewidth}
    \textbf{a)} $x + 51\,340 = 98\,210$\par
    \dongke{3}
    \vspace{0.4mm}
    \textbf{c)} $96\,400 - x = 52\,730$\par
    \dongke{3}
  \end{minipage}\hfill
  \begin{minipage}[t]{0.485\linewidth}
    \textbf{b)} $x - 62\,450 = 83\,710$\par
    \dongke{3}
    \vspace{0.4mm}
    \textbf{d)} $x - (24\,500 + 15\,500) = 63\,000$\par
    \dongke{3}
  \end{minipage}
}

\vspace{0.8mm}

% BÀI 3: ĐỔI ĐƠN VỊ (4 Ý - 1 CỘT TỪ TRÊN XUỐNG)
\khungBaiTap{%
  \textbf{Bài 3 (1,0 điểm). Viết số thích hợp vào chỗ chấm:}\par\vspace{0.8mm}
  \noindent
  a) $6\,\text{tấn}\,540\,\text{kg} = \dotfill\,\text{kg}$\par\vspace{1.4mm}
  \noindent
  b) $9\,\text{m}^2\,65\,\text{dm}^2 = \dotfill\,\text{dm}^2$\par\vspace{1.4mm}
  \noindent
  c) $6\,\text{thế kỉ}\,20\,\text{năm} = \dotfill\,\text{năm}$\par\vspace{1.4mm}
  \noindent
  d) $5\,\text{phút}\,18\,\text{giây} = \dotfill\,\text{giây}$\par
}

\vspace{0.8mm}

% BÀI 4: BÀI TOÁN GIẢI (11 DÒNG KẺ)
\khungBaiTap{%
  \textbf{Bài 4 (1,5 điểm).} Một xưởng may nhận may 150 bộ quần áo. Ngày đầu xưởng may được 30\% số bộ quần áo. Ngày thứ hai may được 45\% số bộ quần áo còn lại. Mỗi bộ quần áo xưởng nhận được 120 000 đồng tiền công. Hỏi sau hai ngày, xưởng thu được bao nhiêu tiền công?
  \par\vspace{0.3mm}%
  \noindent\textit{Bài giải:}
  \dongke{11}
}

\vspace{0.8mm}

% BÀI 5: TOÁN NÂNG CAO (8 DÒNG KẺ)
\khungBaiTap{%
  \textbf{Bài 5 (1,0 điểm).} Tính bằng cách thuận tiện nhất:\par\vspace{0.6mm}\centerline{$K = 0{,}25 \times 4 \times 6{,}7$}
  \par\vspace{0.3mm}%
  \noindent\textit{Bài làm:}
  \dongke{8}
}

\newpage


% =============================================================================
% ========================= MÃ ĐỀ: 404 ========================================
% =============================================================================

% -------------------- TRANG 1: MẶT TRƯỚC --------------------
\noindent
\begin{tabularx}{\linewidth}{|X|c|}
\hline
\textbf{PHIẾU BÀI TẬP CUỐI TUẦN -- TOÁN LỚP 4 (TUẦN 12)} & \textbf{MÃ ĐỀ: 404} \\
\textit{Trung tâm Toán Thầy Lĩnh -- Năm học: 2026 - 2027} & \textit{Thời gian: 45 phút} \\
\hline
\end{tabularx}

\vspace{1mm}

\noindent
\begin{tabularx}{\linewidth}{|X|p{3.5cm}|}
\hline
\rule[-2mm]{0pt}{7.5mm}\textbf{Họ và tên:} \dotfill & \textbf{Lớp:} \makebox[2cm]{\dotfill} \\
\hline
\end{tabularx}

\vspace{1mm}

\noindent
\begin{tabularx}{\linewidth}{|p{3cm}|X|}
\hline
\centering\textbf{ĐIỂM SỐ} & \centering\textbf{LỜI NHẬN XÉT CỦA THẦY CÔ} \tabularnewline
\hline
\rule{0pt}{1.1cm} & \tabularnewline
\hline
\end{tabularx}

\vspace{1mm}

\noindent\textbf{BẢNG ĐIỀN ĐÁP ÁN TRẮC NGHIỆM (Mỗi câu đúng 0,25 điểm):}
\vspace{0.6mm}

\begin{tabularx}{\linewidth}{|c|*{12}{>{\centering\arraybackslash}X|}}
\hline
\textbf{Câu} & \textbf{1} & \textbf{2} & \textbf{3} & \textbf{4} & \textbf{5} & \textbf{6} & \textbf{7} & \textbf{8} & \textbf{9} & \textbf{10} & \textbf{11} & \textbf{12} \\
\hline
\textbf{Đáp án} & & & & & & & & & & & & \\
\hline
\end{tabularx}

\vspace{1.2mm}

\noindent\textbf{\large PHẦN I. TRẮC NGHIỆM KHÁCH QUAN (3,0 điểm)}

\cauTN{1}{Kết quả của phép tính $924\,500 - 568\,342$ là:}{$356\,158$}{$356\,258$}{$366\,158$}{$355\,158$}
\cauTN{2}{Điền số thích hợp: $(46 + 89) + 11 = 46 + (89 + \dots)$}{$11$}{$46$}{$89$}{$100$}
\cauTN{3}{Tính thuận tiện: $250 + 380 + 750 = \dots$}{$1\,280$}{$1\,380$}{$1\,480$}{$1\,350$}
\cauTN{4}{Số thích hợp điền vào chỗ chấm $4\,\text{tấn}\,780\,\text{kg} = \dots\,\text{kg}$ là:}{$478$}{$4\,780$}{$47\,800$}{$40\,780$}
\cauTN{5}{Số thích hợp điền vào chỗ chấm $5\,\text{m}^2\,84\,\text{dm}^2 = \dots\,\text{dm}^2$ là:}{$584$}{$5\,084$}{$5\,840$}{$59$}
\cauTN{6}{Khoảng thời gian $3\,\text{phút}\,55\,\text{giây}$ bằng bao nhiêu giây?}{$225\,\text{giây}$}{$235\,\text{giây}$}{$245\,\text{giây}$}{$255\,\text{giây}$}
\cauTN{7}{Tìm $x$: $x - 56\,240 = 78\,390$:}{$x = 134\,630$}{$x = 22\,150$}{$x = 134\,530$}{$x = 135\,630$}
\cauTN{8}{Điền dấu thích hợp: $78\,000 - 32\,000 \;\dots\; 80\,000 - 35\,000$}{$>$}{$<$}{$=$}{Không so sánh được}
\cauTN{9}{Hiệu của số tròn chục nghìn lớn nhất có 5 chữ số và số tròn chục nghìn bé nhất có 5 chữ số là:}{$80\,000$}{$90\,000$}{$89\,000$}{$70\,000$}
\cauTN{10}{Làm tròn số $956\,200$ đến hàng chục nghìn ta được số:}{$950\,000$}{$960\,000$}{$1\,000\,000$}{$956\,000$}
\cauTNhaiDong{11}{Chu vi hình chữ nhật có chiều dài $36\,\text{m}$, chiều rộng $24\,\text{m}$ là:}{$60\,\text{m}$}{$120\,\text{m}$}{$864\,\text{m}^2$}{$140\,\text{m}$}
\cauTNhaiDong{12}{Một bể bơi chứa $60\,000\,\text{lít}$ nước, người ta đã xả ra $25\,000\,\text{lít}$. Trong bể còn lại số lít nước là:}{$35\,000\,\text{lít}$}{$45\,000\,\text{lít}$}{$30\,000\,\text{lít}$}{$25\,000\,\text{lít}$}

\vspace{0.4mm}
\noindent\textbf{\large PHẦN II. TỰ LUẬN (7,0 điểm)}

\noindent\textbf{Bài 1 (2,0 điểm). Đặt tính rồi tính:}
\vspace{0.4mm}

\noindent
\begin{tabularx}{\linewidth}{|X|X|X|X|}
\hline
\oDatTinh[4]{a)}{$852\,340 - 486\,579$} & 
\oDatTinh[4]{b)}{$800\,000 - 347\,628$} & 
\oDatTinh[4]{c)}{$537\,486 + 284\,635$} & 
\oDatTinh[4]{d)}{$682\,145 + 217\,855$} \tabularnewline
\hline
\end{tabularx}

\newpage
% -------------------- TRANG 2: MẶT SAU --------------------
\noindent
\begin{tabularx}{\linewidth}{|X|c|}
\hline
\textbf{PHẦN TỰ LUẬN (TIẾP THEO) -- PHIẾU BÀI TẬP TUẦN 12} & \textbf{MÃ ĐỀ: 404} \\
\hline
\end{tabularx}

\vspace{0.8mm}

% BÀI 2: TÌM X (4 Ý)
\khungBaiTap{%
  \textbf{Bài 2 (1,5 điểm). Tìm $x$:}\par\vspace{0.6mm}
  \noindent
  \begin{minipage}[t]{0.485\linewidth}
    \textbf{a)} $x + 48\,620 = 94\,150$\par
    \dongke{3}
    \vspace{0.4mm}
    \textbf{c)} $94\,600 - x = 48\,290$\par
    \dongke{3}
  \end{minipage}\hfill
  \begin{minipage}[t]{0.485\linewidth}
    \textbf{b)} $x - 57\,380 = 76\,450$\par
    \dongke{3}
    \vspace{0.4mm}
    \textbf{d)} $x - (28\,300 + 11\,700) = 55\,000$\par
    \dongke{3}
  \end{minipage}
}

\vspace{0.8mm}

% BÀI 3: ĐỔI ĐƠN VỊ (4 Ý - 1 CỘT TỪ TRÊN XUỐNG)
\khungBaiTap{%
  \textbf{Bài 3 (1,0 điểm). Viết số thích hợp vào chỗ chấm:}\par\vspace{0.8mm}
  \noindent
  a) $5\,\text{tấn}\,850\,\text{kg} = \dotfill\,\text{kg}$\par\vspace{1.4mm}
  \noindent
  b) $8\,\text{m}^2\,94\,\text{dm}^2 = \dotfill\,\text{dm}^2$\par\vspace{1.4mm}
  \noindent
  c) $7\,\text{thế kỉ}\,40\,\text{năm} = \dotfill\,\text{năm}$\par\vspace{1.4mm}
  \noindent
  d) $4\,\text{phút}\,40\,\text{giây} = \dotfill\,\text{giây}$\par
}

\vspace{0.8mm}

% BÀI 4: BÀI TOÁN GIẢI (11 DÒNG KẺ)
\khungBaiTap{%
  \textbf{Bài 4 (1,5 điểm).} Một cửa hàng bán 3{,}6 tấn gạo trong ba ngày. Ngày thứ nhất bán được 25\% số gạo. Ngày thứ hai bán được 60\% số gạo còn lại. Mỗi tấn gạo cửa hàng bán được lãi 1 200 000 đồng. Hỏi ngày thứ ba cửa hàng bán được bao nhiêu tấn gạo và tiền lãi bán gạo ngày thứ ba là bao nhiêu?
  \par\vspace{0.3mm}%
  \noindent\textit{Bài giải:}
  \dongke{11}
}

\vspace{0.8mm}

% BÀI 5: TOÁN NÂNG CAO (8 DÒNG KẺ)
\khungBaiTap{%
  \textbf{Bài 5 (1,0 điểm).} Tính bằng cách thuận tiện nhất:\par\vspace{0.6mm}\centerline{$K = 15{,}4 \times 101 - 15{,}4$}
  \par\vspace{0.3mm}%
  \noindent\textit{Bài làm:}
  \dongke{8}
}

\newpage


% =============================================================================
% ========================= MÃ ĐỀ: 405 ========================================
% =============================================================================

% -------------------- TRANG 1: MẶT TRƯỚC --------------------
\noindent
\begin{tabularx}{\linewidth}{|X|c|}
\hline
\textbf{PHIẾU BÀI TẬP CUỐI TUẦN -- TOÁN LỚP 4 (TUẦN 12)} & \textbf{MÃ ĐỀ: 405} \\
\textit{Trung tâm Toán Thầy Lĩnh -- Năm học: 2026 - 2027} & \textit{Thời gian: 45 phút} \\
\hline
\end{tabularx}

\vspace{1mm}

\noindent
\begin{tabularx}{\linewidth}{|X|p{3.5cm}|}
\hline
\rule[-2mm]{0pt}{7.5mm}\textbf{Họ và tên:} \dotfill & \textbf{Lớp:} \makebox[2cm]{\dotfill} \\
\hline
\end{tabularx}

\vspace{1mm}

\noindent
\begin{tabularx}{\linewidth}{|p{3cm}|X|}
\hline
\centering\textbf{ĐIỂM SỐ} & \centering\textbf{LỜI NHẬN XÉT CỦA THẦY CÔ} \tabularnewline
\hline
\rule{0pt}{1.1cm} & \tabularnewline
\hline
\end{tabularx}

\vspace{1mm}

\noindent\textbf{BẢNG ĐIỀN ĐÁP ÁN TRẮC NGHIỆM (Mỗi câu đúng 0,25 điểm):}
\vspace{0.6mm}

\begin{tabularx}{\linewidth}{|c|*{12}{>{\centering\arraybackslash}X|}}
\hline
\textbf{Câu} & \textbf{1} & \textbf{2} & \textbf{3} & \textbf{4} & \textbf{5} & \textbf{6} & \textbf{7} & \textbf{8} & \textbf{9} & \textbf{10} & \textbf{11} & \textbf{12} \\
\hline
\textbf{Đáp án} & & & & & & & & & & & & \\
\hline
\end{tabularx}

\vspace{1.2mm}

\noindent\textbf{\large PHẦN I. TRẮC NGHIỆM KHÁCH QUAN (3,0 điểm)}

\cauTN{1}{Kết quả của phép tính $765\,410 - 389\,654$ là:}{$375\,756$}{$375\,856$}{$385\,756$}{$374\,756$}
\cauTN{2}{Muốn cộng một tổng hai số với số thứ ba, ta có thể:}{Cộng số thứ nhất với tổng của số thứ hai và số thứ ba}{Lấy số thứ nhất nhân số thứ hai}{Lấy số thứ ba trừ đi tổng}{Nhân ba số với nhau}
\cauTN{3}{Tính thuận tiện: $460 + 275 + 140 = \dots$}{$865$}{$875$}{$885$}{$855$}
\cauTN{4}{Số thích hợp điền vào chỗ chấm $7\,\text{tấn}\,920\,\text{kg} = \dots\,\text{kg}$ là:}{$792$}{$7\,920$}{$79\,200$}{$70\,920$}
\cauTN{5}{Số thích hợp điền vào chỗ chấm $6\,\text{m}^2\,47\,\text{dm}^2 = \dots\,\text{dm}^2$ là:}{$647$}{$6\,047$}{$6\,470$}{$65$}
\cauTN{6}{Khoảng thời gian $4\,\text{phút}\,20\,\text{giây}$ bằng bao nhiêu giây?}{$240\,\text{giây}$}{$260\,\text{giây}$}{$280\,\text{giây}$}{$300\,\text{giây}$}
\cauTN{7}{Tìm $x$: $86\,400 - x = 42\,850$:}{$x = 43\,550$}{$x = 43\,650$}{$x = 129\,250$}{$x = 42\,550$}
\cauTN{8}{Điền dấu thích hợp: $84\,000 - 36\,000 \;\dots\; 95\,000 - 45\,000$}{$>$}{$<$}{$=$}{Không so sánh được}
\cauTN{9}{Hiệu giữa số lớn nhất có 5 chữ số khác nhau và số bé nhất có 5 chữ số khác nhau là:}{$88\,531$}{$88\,530$}{$87\,531$}{$89\,531$}
\cauTN{10}{Làm tròn số $674\,900$ đến hàng chục nghìn ta được số:}{$670\,000$}{$680\,000$}{$674\,000$}{$700\,000$}
\cauTNhaiDong{11}{Một hình tam giác có ba cạnh bằng nhau, mỗi cạnh dài $15\,\text{cm}$. Chu vi hình tam giác đó là:}{$30\,\text{cm}$}{$45\,\text{cm}$}{$60\,\text{cm}$}{$50\,\text{cm}$}
\cauTNhaiDong{12}{Một đoàn tàu chở $75\,000\,\text{kg}$ hàng, ga thứ nhất bốc dỡ xuống $28\,000\,\text{kg}$. Trên tàu còn lại số tấn hàng là:}{$47\,\text{tấn}$}{$48\,\text{tấn}$}{$45\,\text{tấn}$}{$50\,\text{tấn}$}

\vspace{0.4mm}
\noindent\textbf{\large PHẦN II. TỰ LUẬN (7,0 điểm)}

\noindent\textbf{Bài 1 (2,0 điểm). Đặt tính rồi tính:}
\vspace{0.4mm}

\noindent
\begin{tabularx}{\linewidth}{|X|X|X|X|}
\hline
\oDatTinh[4]{a)}{$784\,520 - 396\,785$} & 
\oDatTinh[4]{b)}{$950\,000 - 482\,635$} & 
\oDatTinh[4]{c)}{$469\,358 + 375\,864$} & 
\oDatTinh[4]{d)}{$543\,210 + 456\,790$} \tabularnewline
\hline
\end{tabularx}

\newpage
% -------------------- TRANG 2: MẶT SAU --------------------
\noindent
\begin{tabularx}{\linewidth}{|X|c|}
\hline
\textbf{PHẦN TỰ LUẬN (TIẾP THEO) -- PHIẾU BÀI TẬP TUẦN 12} & \textbf{MÃ ĐỀ: 405} \\
\hline
\end{tabularx}

\vspace{0.8mm}

% BÀI 2: TÌM X (4 Ý)
\khungBaiTap{%
  \textbf{Bài 2 (1,5 điểm). Tìm $x$:}\par\vspace{0.6mm}
  \noindent
  \begin{minipage}[t]{0.485\linewidth}
    \textbf{a)} $x + 46\,390 = 92\,540$\par
    \dongke{3}
    \vspace{0.4mm}
    \textbf{c)} $93\,500 - x = 47\,280$\par
    \dongke{3}
  \end{minipage}\hfill
  \begin{minipage}[t]{0.485\linewidth}
    \textbf{b)} $x - 48\,620 = 79\,530$\par
    \dongke{3}
    \vspace{0.4mm}
    \textbf{d)} $x - (32\,400 + 17\,600) = 48\,000$\par
    \dongke{3}
  \end{minipage}
}

\vspace{0.8mm}

% BÀI 3: ĐỔI ĐƠN VỊ (4 Ý - 1 CỘT TỪ TRÊN XUỐNG)
\khungBaiTap{%
  \textbf{Bài 3 (1,0 điểm). Viết số thích hợp vào chỗ chấm:}\par\vspace{0.8mm}
  \noindent
  a) $7\,\text{tấn}\,640\,\text{kg} = \dotfill\,\text{kg}$\par\vspace{1.4mm}
  \noindent
  b) $9\,\text{m}^2\,48\,\text{dm}^2 = \dotfill\,\text{dm}^2$\par\vspace{1.4mm}
  \noindent
  c) $8\,\text{thế kỉ}\,15\,\text{năm} = \dotfill\,\text{năm}$\par\vspace{1.4mm}
  \noindent
  d) $3\,\text{phút}\,45\,\text{giây} = \dotfill\,\text{giây}$\par
}

\vspace{0.8mm}

% BÀI 4: BÀI TOÁN GIẢI (11 DÒNG KẺ)
\khungBaiTap{%
  \textbf{Bài 4 (1,5 điểm).} Một mảnh vườn hình chữ nhật có chiều dài 18{,}5 m và chiều rộng 12 m. Người ta trồng rau trên 70\% diện tích mảnh vườn, phần còn lại làm lối đi. Cứ mỗi mét vuông trồng rau thu hoạch được 2{,}5 kg rau, giá bán 1 kg rau là 12 000 đồng. Hỏi số rau thu hoạch được bán hết thì thu được bao nhiêu tiền?
  \par\vspace{0.3mm}%
  \noindent\textit{Bài giải:}
  \dongke{11}
}

\vspace{0.8mm}

% BÀI 5: TOÁN NÂNG CAO (8 DÒNG KẺ)
\khungBaiTap{%
  \textbf{Bài 5 (1,0 điểm).} Tính bằng cách thuận tiện nhất:\par\vspace{0.6mm}\centerline{$K = 7{,}6 \times 9{,}9 + 7{,}6 \times 0{,}1$}
  \par\vspace{0.3mm}%
  \noindent\textit{Bài làm:}
  \dongke{8}
}

\newpage


% =============================================================================
% ========================= MÃ ĐỀ: 406 ========================================
% =============================================================================

% -------------------- TRANG 1: MẶT TRƯỚC --------------------
\noindent
\begin{tabularx}{\linewidth}{|X|c|}
\hline
\textbf{PHIẾU BÀI TẬP CUỐI TUẦN -- TOÁN LỚP 4 (TUẦN 12)} & \textbf{MÃ ĐỀ: 406} \\
\textit{Trung tâm Toán Thầy Lĩnh -- Năm học: 2026 - 2027} & \textit{Thời gian: 45 phút} \\
\hline
\end{tabularx}

\vspace{1mm}

\noindent
\begin{tabularx}{\linewidth}{|X|p{3.5cm}|}
\hline
\rule[-2mm]{0pt}{7.5mm}\textbf{Họ và tên:} \dotfill & \textbf{Lớp:} \makebox[2cm]{\dotfill} \\
\hline
\end{tabularx}

\vspace{1mm}

\noindent
\begin{tabularx}{\linewidth}{|p{3cm}|X|}
\hline
\centering\textbf{ĐIỂM SỐ} & \centering\textbf{LỜI NHẬN XÉT CỦA THẦY CÔ} \tabularnewline
\hline
\rule{0pt}{1.1cm} & \tabularnewline
\hline
\end{tabularx}

\vspace{1mm}

\noindent\textbf{BẢNG ĐIỀN ĐÁP ÁN TRẮC NGHIỆM (Mỗi câu đúng 0,25 điểm):}
\vspace{0.6mm}

\begin{tabularx}{\linewidth}{|c|*{12}{>{\centering\arraybackslash}X|}}
\hline
\textbf{Câu} & \textbf{1} & \textbf{2} & \textbf{3} & \textbf{4} & \textbf{5} & \textbf{6} & \textbf{7} & \textbf{8} & \textbf{9} & \textbf{10} & \textbf{11} & \textbf{12} \\
\hline
\textbf{Đáp án} & & & & & & & & & & & & \\
\hline
\end{tabularx}

\vspace{1.2mm}

\noindent\textbf{\large PHẦN I. TRẮC NGHIỆM KHÁCH QUAN (3,0 điểm)}

\cauTN{1}{Kết quả của phép tính $834\,620 - 467\,855$ là:}{$366\,765$}{$366\,865$}{$376\,765$}{$365\,765$}
\cauTN{2}{Công thức biểu diễn tính chất kết hợp của phép cộng là:}{$(a+b)+c = a+(b+c)$}{$a+b=b+a$}{$a-b=b-a$}{$a \times b = b \times a$}
\cauTN{3}{Tính thuận tiện: $380 + 145 + 120 = \dots$}{$635$}{$645$}{$655$}{$625$}
\cauTN{4}{Số thích hợp điền vào chỗ chấm $3\,\text{tấn}\,580\,\text{kg} = \dots\,\text{kg}$ là:}{$358$}{$3\,580$}{$35\,800$}{$30\,580$}
\cauTN{5}{Số thích hợp điền vào chỗ chấm $7\,\text{m}^2\,63\,\text{dm}^2 = \dots\,\text{dm}^2$ là:}{$763$}{$7\,063$}{$7\,630$}{$79$}
\cauTN{6}{Khoảng thời gian $5\,\text{phút}\,10\,\text{giây}$ bằng bao nhiêu giây?}{$300\,\text{giây}$}{$310\,\text{giây}$}{$320\,\text{giây}$}{$330\,\text{giây}$}
\cauTN{7}{Tìm $x$: $91\,500 - x = 45\,380$:}{$x = 46\,120$}{$x = 46\,220$}{$x = 136\,880$}{$x = 45\,120$}
\cauTN{8}{Điền dấu thích hợp: $90\,000 - 38\,000 \;\dots\; 85\,000 - 32\,000$}{$>$}{$<$}{$=$}{Không so sánh được}
\cauTN{9}{Hiệu giữa $100\,000$ và số liền trước nó là:}{$0$}{$1$}{$2$}{$10$}
\cauTN{10}{Làm tròn số $785\,600$ đến hàng chục nghìn ta được số:}{$780\,000$}{$790\,000$}{$800\,000$}{$785\,000$}
\cauTNhaiDong{11}{Một mảnh đất hình tam giác có ba cạnh là $28\,\text{m}, 32\,\text{m}, 40\,\text{m}$. Chu vi mảnh đất là:}{$90\,\text{m}$}{$100\,\text{m}$}{$110\,\text{m}$}{$120\,\text{m}$}
\cauTNhaiDong{12}{Một nông trường có $80\,000\,\text{cây}$ ăn quả, vừa qua đã thu hoạch xong $35\,000\,\text{cây}$. Số cây chưa thu hoạch là:}{$45\,000\,\text{cây}$}{$55\,000\,\text{cây}$}{$40\,000\,\text{cây}$}{$50\,000\,\text{cây}$}

\vspace{0.4mm}
\noindent\textbf{\large PHẦN II. TỰ LUẬN (7,0 điểm)}

\noindent\textbf{Bài 1 (2,0 điểm). Đặt tính rồi tính:}
\vspace{0.4mm}

\noindent
\begin{tabularx}{\linewidth}{|X|X|X|X|}
\hline
\oDatTinh[4]{a)}{$846\,210 - 478\,565$} & 
\oDatTinh[4]{b)}{$920\,000 - 536\,428$} & 
\oDatTinh[4]{c)}{$527\,648 + 286\,759$} & 
\oDatTinh[4]{d)}{$625\,340 + 374\,660$} \tabularnewline
\hline
\end{tabularx}

\newpage
% -------------------- TRANG 2: MẶT SAU --------------------
\noindent
\begin{tabularx}{\linewidth}{|X|c|}
\hline
\textbf{PHẦN TỰ LUẬN (TIẾP THEO) -- PHIẾU BÀI TẬP TUẦN 12} & \textbf{MÃ ĐỀ: 406} \\
\hline
\end{tabularx}

\vspace{0.8mm}

% BÀI 2: TÌM X (4 Ý)
\khungBaiTap{%
  \textbf{Bài 2 (1,5 điểm). Tìm $x$:}\par\vspace{0.6mm}
  \noindent
  \begin{minipage}[t]{0.485\linewidth}
    \textbf{a)} $x + 53\,480 = 96\,210$\par
    \dongke{3}
    \vspace{0.4mm}
    \textbf{c)} $97\,300 - x = 54\,620$\par
    \dongke{3}
  \end{minipage}\hfill
  \begin{minipage}[t]{0.485\linewidth}
    \textbf{b)} $x - 52\,840 = 84\,670$\par
    \dongke{3}
    \vspace{0.4mm}
    \textbf{d)} $x - (35\,200 + 14\,800) = 46\,000$\par
    \dongke{3}
  \end{minipage}
}

\vspace{0.8mm}

% BÀI 3: ĐỔI ĐƠN VỊ (4 Ý - 1 CỘT TỪ TRÊN XUỐNG)
\khungBaiTap{%
  \textbf{Bài 3 (1,0 điểm). Viết số thích hợp vào chỗ chấm:}\par\vspace{0.8mm}
  \noindent
  a) $8\,\text{tấn}\,420\,\text{kg} = \dotfill\,\text{kg}$\par\vspace{1.4mm}
  \noindent
  b) $6\,\text{m}^2\,85\,\text{dm}^2 = \dotfill\,\text{dm}^2$\par\vspace{1.4mm}
  \noindent
  c) $5\,\text{thế kỉ}\,45\,\text{năm} = \dotfill\,\text{năm}$\par\vspace{1.4mm}
  \noindent
  d) $4\,\text{phút}\,55\,\text{giây} = \dotfill\,\text{giây}$\par
}

\vspace{0.8mm}

% BÀI 4: BÀI TOÁN GIẢI (11 DÒNG KẺ)
\khungBaiTap{%
  \textbf{Bài 4 (1,5 điểm).} Một cửa hàng nhập 2{,}4 tấn đường. Ngày thứ nhất bán được 35\% số đường, ngày thứ hai bán được 50\% số đường còn lại. Mỗi ki-lô-gam đường bán ra cửa hàng lãi 2 500 đồng. Hỏi sau hai ngày, cửa hàng lãi được bao nhiêu tiền từ số đường đã bán?
  \par\vspace{0.3mm}%
  \noindent\textit{Bài giải:}
  \dongke{11}
}

\vspace{0.8mm}

% BÀI 5: TOÁN NÂNG CAO (8 DÒNG KẺ)
\khungBaiTap{%
  \textbf{Bài 5 (1,0 điểm).} Tính bằng cách thuận tiện nhất:\par\vspace{0.6mm}\centerline{$K = 0{,}25 \times 4 \times 12{,}8$}
  \par\vspace{0.3mm}%
  \noindent\textit{Bài làm:}
  \dongke{8}
}

\newpage


% =============================================================================
% ========================= MÃ ĐỀ: 407 ========================================
% =============================================================================

% -------------------- TRANG 1: MẶT TRƯỚC --------------------
\noindent
\begin{tabularx}{\linewidth}{|X|c|}
\hline
\textbf{PHIẾU BÀI TẬP CUỐI TUẦN -- TOÁN LỚP 4 (TUẦN 12)} & \textbf{MÃ ĐỀ: 407} \\
\textit{Trung tâm Toán Thầy Lĩnh -- Năm học: 2026 - 2027} & \textit{Thời gian: 45 phút} \\
\hline
\end{tabularx}

\vspace{1mm}

\noindent
\begin{tabularx}{\linewidth}{|X|p{3.5cm}|}
\hline
\rule[-2mm]{0pt}{7.5mm}\textbf{Họ và tên:} \dotfill & \textbf{Lớp:} \makebox[2cm]{\dotfill} \\
\hline
\end{tabularx}

\vspace{1mm}

\noindent
\begin{tabularx}{\linewidth}{|p{3cm}|X|}
\hline
\centering\textbf{ĐIỂM SỐ} & \centering\textbf{LỜI NHẬN XÉT CỦA THẦY CÔ} \tabularnewline
\hline
\rule{0pt}{1.1cm} & \tabularnewline
\hline
\end{tabularx}

\vspace{1mm}

\noindent\textbf{BẢNG ĐIỀN ĐÁP ÁN TRẮC NGHIỆM (Mỗi câu đúng 0,25 điểm):}
\vspace{0.6mm}

\begin{tabularx}{\linewidth}{|c|*{12}{>{\centering\arraybackslash}X|}}
\hline
\textbf{Câu} & \textbf{1} & \textbf{2} & \textbf{3} & \textbf{4} & \textbf{5} & \textbf{6} & \textbf{7} & \textbf{8} & \textbf{9} & \textbf{10} & \textbf{11} & \textbf{12} \\
\hline
\textbf{Đáp án} & & & & & & & & & & & & \\
\hline
\end{tabularx}

\vspace{1.2mm}

\noindent\textbf{\large PHẦN I. TRẮC NGHIỆM KHÁCH QUAN (3,0 điểm)}

\cauTN{1}{Kết quả của phép tính $915\,300 - 548\,675$ là:}{$366\,625$}{$366\,725$}{$376\,625$}{$365\,625$}
\cauTN{2}{Phép tính $125 + 75 + 340 = 200 + 340$ là áp dụng tính chất:}{Giao hoán}{Kết hợp}{Phân phối}{Không áp dụng}
\cauTN{3}{Tính thuận tiện: $520 + 195 + 180 = \dots$}{$885$}{$895$}{$905$}{$875$}
\cauTN{4}{Số thích hợp điền vào chỗ chấm $9\,\text{tấn}\,240\,\text{kg} = \dots\,\text{kg}$ là:}{$924$}{$9\,240$}{$92\,400$}{$90\,240$}
\cauTN{5}{Số thích hợp điền vào chỗ chấm $8\,\text{m}^2\,37\,\text{dm}^2 = \dots\,\text{dm}^2$ là:}{$837$}{$8\,037$}{$8\,370$}{$89$}
\cauTN{6}{Khoảng thời gian $6\,\text{phút}\,15\,\text{giây}$ bằng bao nhiêu giây?}{$360\,\text{giây}$}{$375\,\text{giây}$}{$385\,\text{giây}$}{$395\,\text{giây}$}
\cauTN{7}{Tìm $x$: $95\,600 - x = 51\,420$:}{$x = 44\,180$}{$x = 44\,280$}{$x = 147\,020$}{$x = 43\,180$}
\cauTN{8}{Điền dấu thích hợp: $95\,000 - 42\,000 \;\dots\; 88\,000 - 35\,000$}{$>$}{$<$}{$=$}{Không so sánh được}
\cauTN{9}{Hiệu giữa số lớn nhất có 6 chữ số và số tròn trăm nghìn lớn nhất có 6 chữ số là:}{$99\,999$}{$100\,000$}{$9\,999$}{$90\,000$}
\cauTN{10}{Làm tròn số $894\,200$ đến hàng chục nghìn ta được số:}{$890\,000$}{$900\,000$}{$895\,000$}{$880\,000$}
\cauTNhaiDong{11}{Chu vi một hình vuông có cạnh $35\,\text{m}$ là:}{$70\,\text{m}$}{$140\,\text{m}$}{$1\,225\,\text{m}^2$}{$120\,\text{m}$}
\cauTNhaiDong{12}{Một đầm tôm đợt 1 thu hoạch được $12\,500\,\text{kg}$ tôm, đợt 2 thu hoạch được $8\,500\,\text{kg}$. Tổng sản lượng cả hai đợt tính theo tấn là:}{$20\,\text{tấn}$}{$21\,\text{tấn}$}{$22\,\text{tấn}$}{$210\,\text{tấn}$}

\vspace{0.4mm}
\noindent\textbf{\large PHẦN II. TỰ LUẬN (7,0 điểm)}

\noindent\textbf{Bài 1 (2,0 điểm). Đặt tính rồi tính:}
\vspace{0.4mm}

\noindent
\begin{tabularx}{\linewidth}{|X|X|X|X|}
\hline
\oDatTinh[4]{a)}{$865\,420 - 489\,675$} & 
\oDatTinh[4]{b)}{$910\,000 - 473\,582$} & 
\oDatTinh[4]{c)}{$476\,539 + 385\,684$} & 
\oDatTinh[4]{d)}{$748\,125 + 251\,875$} \tabularnewline
\hline
\end{tabularx}

\newpage
% -------------------- TRANG 2: MẶT SAU --------------------
\noindent
\begin{tabularx}{\linewidth}{|X|c|}
\hline
\textbf{PHẦN TỰ LUẬN (TIẾP THEO) -- PHIẾU BÀI TẬP TUẦN 12} & \textbf{MÃ ĐỀ: 407} \\
\hline
\end{tabularx}

\vspace{0.8mm}

% BÀI 2: TÌM X (4 Ý)
\khungBaiTap{%
  \textbf{Bài 2 (1,5 điểm). Tìm $x$:}\par\vspace{0.6mm}
  \noindent
  \begin{minipage}[t]{0.485\linewidth}
    \textbf{a)} $x + 56\,340 = 98\,120$\par
    \dongke{3}
    \vspace{0.4mm}
    \textbf{c)} $98\,400 - x = 56\,370$\par
    \dongke{3}
  \end{minipage}\hfill
  \begin{minipage}[t]{0.485\linewidth}
    \textbf{b)} $x - 63\,520 = 87\,490$\par
    \dongke{3}
    \vspace{0.4mm}
    \textbf{d)} $x - (41\,500 + 18\,500) = 54\,000$\par
    \dongke{3}
  \end{minipage}
}

\vspace{0.8mm}

% BÀI 3: ĐỔI ĐƠN VỊ (4 Ý - 1 CỘT TỪ TRÊN XUỐNG)
\khungBaiTap{%
  \textbf{Bài 3 (1,0 điểm). Viết số thích hợp vào chỗ chấm:}\par\vspace{0.8mm}
  \noindent
  a) $7\,\text{tấn}\,750\,\text{kg} = \dotfill\,\text{kg}$\par\vspace{1.4mm}
  \noindent
  b) $8\,\text{m}^2\,52\,\text{dm}^2 = \dotfill\,\text{dm}^2$\par\vspace{1.4mm}
  \noindent
  c) $6\,\text{thế kỉ}\,50\,\text{năm} = \dotfill\,\text{năm}$\par\vspace{1.4mm}
  \noindent
  d) $5\,\text{phút}\,40\,\text{giây} = \dotfill\,\text{giây}$\par
}

\vspace{0.8mm}

% BÀI 4: BÀI TOÁN GIẢI (11 DÒNG KẺ)
\khungBaiTap{%
  \textbf{Bài 4 (1,5 điểm).} Một gia đình dùng hết 28{,}5 $\text{m}^3$ nước trong tháng. Giá 1 $\text{m}^3$ nước là 9 500 đồng. Ngoài ra, gia đình phải trả thêm 10\% phí bảo vệ môi trường trên tổng tiền nước. Hỏi tổng số tiền gia đình phải trả trong tháng là bao nhiêu?
  \par\vspace{0.3mm}%
  \noindent\textit{Bài giải:}
  \dongke{11}
}

\vspace{0.8mm}

% BÀI 5: TOÁN NÂNG CAO (8 DÒNG KẺ)
\khungBaiTap{%
  \textbf{Bài 5 (1,0 điểm).} Tính bằng cách thuận tiện nhất:\par\vspace{0.6mm}\centerline{$K = 24{,}6 \times 101 - 24{,}6$}
  \par\vspace{0.3mm}%
  \noindent\textit{Bài làm:}
  \dongke{8}
}

\newpage
\end{document}
